\SourceLeaf{58}
\noindent Generalibus via etiam hic iam strata est. Si enim ab utraque parte, quae in
$\sin(h+1)\varphi$ ducta reperiuntur, aequalia, uti fas est, ponantur, eveniet:

\[
\frac{\nu nA^{(h)}}{2}-\frac{\nu nA^{(h+2)}}{2}=(h+1)A^{(h+1)}+\frac{nhA^{(h)}}{2}+\frac{n(h+2)A^{(h+2)}}{2}.
\]

\noindent Soluta aequatione habebitur

\[
A^{(h+2)}=\frac{n(\nu-h)A^{(h)}-2(h+1)A^{(h+1)}}{n(\nu+h+2)},
\]

\noindent qua formula generali, excepto, quod veniret pro $h=0$, $A^{(2)}$, antea iam singulari
disquisitione inventa, continetur lex cuiuslibet coefficientis in serie nostra fictitia,
quatenus nimirum ille per binos antecedentes facili eruitur calculo. Sic, ut exemplum
adsit

\[
A^{(2)}=\frac{2(\nu nA-A^{(1)})}{n(\nu+2)};\qquad A^{(3)}=\frac{n(\nu-1)A^{(1)}-4A^{(2)}}{n(\nu+3)};
\]

\[
A^{(4)}=\frac{n(\nu-2)A^{(2)}-6A^{(3)}}{n(\nu+4)};\quad\mbox{etc.}
\]

$\delta$) Fieri autem potest, existente $\nu$ numero aliquo negativo, ut coefficiens
aliquis, evanescente denominatore, non definiatur; tunc ad series infinitas ante
exhibitas recurrendum est. Idem hoc fieri oportet in coefficientibus binis prioribus $A$
et $A^{(1)}$, quos ut cognitos scala recursionis iamiam supponit. Potest vero singulari
artificio $A^{(1)}$ per $A$ exhiberi. Erat enim

\[
A=1+.\,.\,.\,.+\frac{1}{2^{2h}}{}^{\nu}\mathfrak{B}^{(2h)}.{}^{2h}\mathfrak{B}^{(h)}n^{2h}+.\,.\,.,
\]

\[
A^{(1)}={}^{\nu}\mathfrak{B}^{(1)}n+.\,.\,.+\frac{1}{2^{2h}}{}^{\nu}\mathfrak{B}^{(2h+1)}.{}^{2h+1}\mathfrak{B}^{(h+1)}n^{2h+1}+.\,.\,..
\]

\noindent Est, si sumamus $2nA+A^{(1)}$, huius seriei terminus initialis

\[
=({}^{\nu}\mathfrak{B}^{(1)}+2)n=(\nu+2)n,
\]

\noindent generalis $h^{\text{tus}}$

\[
2n . \frac{1}{2^{2h}}{}^{\nu}\mathfrak{B}^{(2h)}.{}^{2h}\mathfrak{B}^{(h)}n^{2h}+\frac{1}{2^{2h}}{}^{\nu}\mathfrak{B}^{(2h+1)}.{}^{2h+1}\mathfrak{B}^{(h+1)}n^{2h+1},
\]

\noindent qui, cum sit

\[
{}^{\nu}\mathfrak{B}^{(2h+1)}={}^{\nu}\mathfrak{B}^{(2h)} . \frac{\nu-2h}{2h+1}\qquad\mbox{et}\qquad {}^{2h+1}\mathfrak{B}^{(h+1)}={}^{2h}\mathfrak{B}^{(h)} . \frac{2h+1}{h+1},
\]

\noindent sic contrahitur

\[
\frac{1}{2^{2h}}{}^{\nu}\mathfrak{B}^{(2h)}.{}^{2h}\mathfrak{B}^{(h)}\left(2+\frac{\nu-2h}{2h+1} . \frac{2h+1}{h+1}\right)n^{2h+1}
\]

